\documentclass[11pt]{article}
\usepackage[margin=25mm]{geometry}
\usepackage{fontspec,unicode-math}
\setmainfont{Latin Modern Roman}
\setmathfont{Latin Modern Math}
\usepackage{amsmath,amssymb,mathtools}
\usepackage{xurl,hyperref}
\hypersetup{colorlinks=true,urlcolor=blue,linkcolor=blue}
\setcounter{secnumdepth}{0}
\setlength{\parindent}{0pt}
\setlength{\parskip}{5pt}
\setlength{\emergencystretch}{3em}
\newcommand{\sref}[2]{\hyperlink{source:#1:#2}{[S#2]}}
\newcommand{\cataloguescope}[1]{\paragraph{Source question.}#1}
\newcommand{\floor}[1]{\left\lfloor#1\right\rfloor}
\begin{document}
\section{Independence and Density Thresholds for Cuckoo Hashing}\label{statement}

\subsection{Definitions and mathematical statement}

\paragraph{Shared notation.}Standard number systems are denoted by $\mathbb N,\mathbb Z,\mathbb Q,\mathbb R,\mathbb C$; index ranges, computational models and specialized notation follow the cited problem source and the definitions below.


Cuckoo hashing assigns each stored key to one of its hash locations, with at most one key per location. The load is the ratio of keys to locations; $k$-wise independent hash values are jointly uniform for any $k$ distinct keys.
\paragraph{Statement recorded in the supplied source.}
How much universal independence suffices to guarantee that cuckoo hashing works? For the one-table generalization with $d$ hash functions, what is the exact threshold load $n/m$ even when the hash functions are fully random? \sref{statement}{2}
\paragraph{Scope qualification.}
The independence question and load-threshold question are separate. The supplied assessment reports a solution to the fully random threshold question while retaining unresolved hash-family issues; that dated assessment does not turn both questions into unresolved claims.
\paragraph{Source status.}
The supplied catalogue labels this record partially solved. Its source-status metadata was checked against the supplied snapshot on 2026-10-11. This import does not independently establish current open status or mathematical correctness. The catalogue literature assessment is dated 2026-08-21.
\subsection{Short English statement}
Determine how much hash independence is sufficient and find the fully random load threshold for multi-choice cuckoo hashing.
\subsection{Sources}
\begin{itemize}
\item[\hypertarget{source:statement:1}{S1}] Ulam AI and the UnsolvedMath Contributors, UnsolvedMath: A Curated Collection of Open Mathematics Problems, catalogue entry. Catalogue curation and metadata: CC BY 4.0. The supplied source question is retained; definitions and the English summary are editorial additions.
\url{https://huggingface.co/datasets/ulamai/UnsolvedMath}.
\item[\hypertarget{source:statement:2}{S2}] Combinatorics and Probability (2009), Oberwolfach Reports, DOI 10.4171/owr/2009/23, pp. 1225--1302
\url{https://doi.org/10.4171/owr/2009/23}.
\item[\hypertarget{source:statement:3}{S3}] Fountoulakis-Panagiotou sharp cuckoo thresholds
\url{https://arxiv.org/abs/0910.5147}.
\item[\hypertarget{source:statement:4}{S4}] Dietzfelbinger-Schellbach on universal hash classes
\url{https://doi.org/10.1137/1.9781611973068.87}.
\textit{Underlying Oberwolfach Reports material retains its source-specific licence, including CC BY-SA 4.0 where applicable; see THIRD\_PARTY\_NOTICES.md.}
\end{itemize}
\paragraph{Preparation and review.}
Original question attributed to the cited source and UnsolvedMath contributors. Codex (GPT-6; exact serving revision not exposed) assisted with definitions, English summaries, source inspection and formatting. No new human mathematical authorship or independent proof verification is asserted. Catalogue identity checked against the complete supplied registry and existing repository statements; known flagged repetitions and obvious identical conjectures excluded. Full mathematical-equivalence/literature review has not been performed. Statement preparation and final review followed contribution prompts C and D; automated checks verify structure and identity, not open status.
\label{end:statement}
\end{document}
